\documentclass[11pt,a4paper]{article}

\usepackage[margin=1in]{geometry}
\usepackage{amsmath,amssymb,amsthm}
\usepackage{bm}
\usepackage{enumitem}
\usepackage{hyperref}
\usepackage{mathtools}

\setlist[itemize]{itemsep=2pt,topsep=4pt}
\setlist[enumerate]{itemsep=2pt,topsep=4pt}
\setlength{\parskip}{0.5em}
\setlength{\parindent}{0pt}

\newcommand{\vect}[1]{\bm{#1}}
\newcommand{\R}{\mathbb{R}}

\title{AE700 Detailed Study Material\\Lectures 13 to 17}
\author{Prepared from course lectures: Guidance and Control of Unmanned Autonomous Vehicles}
\date{\today}

\begin{document}
\maketitle
\tableofcontents
\newpage

\section{Lecture 13: Reduced-Order Guidance Models and Straight-Line Path Following}

\subsection{Why reduced-order guidance models are needed}
The full 6-DOF rigid-body equations with aerodynamic force and moment models are too complex for real-time outer-loop guidance design. Hence, reduced-order models are used to:
\begin{itemize}
  \item preserve dominant closed-loop behavior of the autopilot-vehicle system,
  \item simplify path-following control design,
  \item avoid repeated use of full aerodynamic force/moment balance equations.
\end{itemize}

\subsection{Autopilot-hold loop abstractions used in guidance}
Common low-order closed-loop models:
\[
\dot V_a=b_{V_a}(V_{a,c}-V_a), \qquad
\dot \phi=b_{\phi}(\phi_c-\phi),
\]
\[
\ddot h=b_{\dot h}(\dot h_c-\dot h)+b_h(h_c-h), \qquad
\ddot \chi=b_{\dot \chi}(\dot \chi_c-\dot \chi)+b_{\chi}(\chi_c-\chi),
\]
\[
\dot \gamma=b_{\gamma}(\gamma_c-\gamma), \qquad
\dot n_{lf}=b_{n_{lf}}(n_{lf,c}-n_{lf}).
\]
All gains are positive and depend on inner-loop implementation.

\subsection{Kinematic model for controlled flight}
Ground-velocity form:
\[
\dot p_n=V_g\cos\chi\cos\gamma,\quad
\dot p_e=V_g\sin\chi\cos\gamma,\quad
\dot h=V_g\sin\gamma.
\]

With airspeed-heading-flight-path variables and wind:
\[
\begin{bmatrix}\dot p_n\\ \dot p_e\\ \dot h\end{bmatrix}
=
V_a\begin{bmatrix}\cos\psi\cos\gamma_a\\ \sin\psi\cos\gamma_a\\ \sin\gamma_a\end{bmatrix}
+
\begin{bmatrix}w_n\\ w_e\\ -w_d\end{bmatrix}.
\]
For constant altitude and negligible vertical wind:
\[
\dot p_n=V_a\cos\psi+w_n,\qquad
\dot p_e=V_a\sin\psi+w_e,\qquad
\dot h=0.
\]

\subsection{Coordinated-turn relation}
For coordinated turning (no lateral body-frame acceleration):
\[
\dot\chi=\frac{g}{V_g}\tan\phi\cos(\chi-\psi).
\]
Using wind-triangle differentiation and constant-airspeed assumption:
\[
\dot\psi=\frac{g}{V_a}\tan\phi.
\]
This highlights a key idea: turning is achieved through bank angle.

\subsection{Accelerated climb and load factor}
Force balance normal to velocity direction (with bank):
\[
F_{lift}\cos\phi = mV_g\dot\gamma + mg\cos\gamma.
\]
Define load factor
\[
n_{lf}\triangleq \frac{F_{lift}}{mg},
\]
then
\[
\dot\gamma=\frac{g}{V_g}(n_{lf}\cos\phi-\cos\gamma).
\]
For constant-climb condition ($\dot\gamma=0$):
\[
n_{lf}=\frac{\cos\gamma}{\cos\phi}.
\]

\subsection{Kinematic guidance model options}
\subsubsection*{Option 1: course and altitude loops}
\[
\dot p_n=V_a\cos\psi+w_n,\quad
\dot p_e=V_a\sin\psi+w_e,
\]
\[
\ddot\chi=b_{\dot\chi}(\dot\chi_c-\dot\chi)+b_{\chi}(\chi_c-\chi),\quad
\ddot h=b_{\dot h}(\dot h_c-\dot h)+b_h(h_c-h),
\]
\[
\dot V_a=b_{V_a}(V_{a,c}-V_a).
\]
Heading from wind triangle (for $\gamma_a=0$):
\[
\psi=\chi-\sin^{-1}\!\left(\frac{-w_n\sin\chi+w_e\cos\chi}{V_a}\right).
\]

\subsubsection*{Option 2: roll as direct input}
\[
\dot p_n=V_a\cos\psi+w_n,\quad
\dot p_e=V_a\sin\psi+w_e,\quad
\dot\psi=\frac{g}{V_a}\tan\phi,
\]
\[
\ddot h=b_{\dot h}(\dot h_c-\dot h)+b_h(h_c-h),\quad
\dot V_a=b_{V_a}(V_{a,c}-V_a),\quad
\dot\phi=b_{\phi}(\phi_c-\phi).
\]

\subsubsection*{Option 3: load-factor command}
\[
\dot p_n=V_a\cos\psi\cos\gamma_a+w_n,\quad
\dot p_e=V_a\sin\psi\cos\gamma_a+w_e,\quad
\dot h=V_a\sin\gamma_a-w_d,
\]
\[
\dot\psi=\frac{g}{V_a}\tan\phi,\quad
\dot\gamma=\frac{g}{V_g}(n_{lf}\cos\phi-\cos\gamma),
\]
\[
\dot n_{lf}=b_n(n_{lf,c}-n_{lf}),\quad
\dot V_a=b_{V_a}(V_{a,c}-V_a),\quad
\dot\phi=b_\phi(\phi_c-\phi).
\]

\subsection{Dynamic point-mass guidance model}
When lift/drag/thrust are retained explicitly:
\[
\dot p_n=V_g\cos\chi\cos\gamma,\quad
\dot p_e=V_g\sin\chi\cos\gamma,\quad
\dot h=V_g\sin\gamma,
\]
\[
\dot V_g=\frac{F_{thrust}}{m}-\frac{F_{drag}}{m}-g\sin\gamma,
\]
\[
\dot\chi=\frac{F_{lift}\sin\phi\cos(\chi-\psi)}{mV_g\cos\gamma},
\quad
\dot\gamma=\frac{F_{lift}}{mV_g}\cos\phi-\frac{g}{V_g}\cos\gamma.
\]
Aerodynamic closure:
\[
L=\frac{1}{2}\rho V_a^2SC_L,\quad
D=\frac{1}{2}\rho V_a^2SC_D,\quad
C_D=C_{D0}+KC_L^2.
\]

\subsection{Trajectory tracking vs path following}
\begin{itemize}
  \item \textbf{Trajectory tracking}: follow a time-parameterized reference.
  \item \textbf{Path following}: converge to geometric path, without prescribed timing.
\end{itemize}
For small UAVs in wind, path following is usually more robust and less control-intensive.

\subsection{Straight-line path geometry}
Desired line:
\[
\mathcal{P}_{line}(r,q)=\{x\in\R^3:\;x=r+\lambda q,\;\lambda\in\R\},
\]
where $r$ is path origin and $q$ is path unit direction. Path-course angle:
\[
\chi_q=\tan^{-1}\!\left(\frac{q_e}{q_n}\right).
\]
Path-frame error:
\[
e_p^P=\begin{bmatrix}e_{px}\\e_{py}\\e_{pz}\end{bmatrix}=R_i^P(p^i-r^i),\quad
R_i^P=
\begin{bmatrix}
\cos\chi_q & \sin\chi_q & 0\\
-\sin\chi_q & \cos\chi_q & 0\\
0 & 0 & 1
\end{bmatrix}.
\]
Cross-track dynamics:
\[
\dot e_{py}=V_g\sin(\chi-\chi_q)\cos\gamma.
\]

\subsection{Longitudinal straight-line tracking}
Desired altitude from projection geometry:
\[
h_d=-r_d+\sqrt{s_n^2+s_e^2}\left(\frac{q_d}{\sqrt{q_n^2+q_e^2}}\right).
\]
Altitude loop:
\[
\ddot h=b_{\dot h}(\dot h_c-\dot h)+b_h(h_c-h).
\]
If $h_c=h_d$ and loop is well tuned, steady-state altitude error to constant/ramp-type commands tends to zero.

\subsection{Lateral straight-line vector-field strategy}
Construct desired relative course:
\[
\chi_d(e_{py})=-\chi_\infty\frac{2}{\pi}\tan^{-1}(k_{path}e_{py}),
\quad
\chi_\infty\in(0,\pi/2],\;k_{path}>0.
\]
Then commanded course:
\[
\chi_c=\chi_q+\chi_d(e_{py}).
\]
Rule-of-thumb:
\[
k_{path}\approx \frac{1}{R_{min}},
\]
with $R_{min}$ as minimum turning radius.

\newpage
\section{Lecture 14: Orbit Following, Carrot-Chasing, and NLGL}

\subsection{Lyapunov proof for straight-line lateral guidance}
Use:
\[
V(e_{py})=\frac{1}{2}e_{py}^2,\qquad
\chi_c=\chi_q-\chi_\infty\frac{2}{\pi}\tan^{-1}(k_{path}e_{py}).
\]
Assuming inner loop makes $\chi\approx\chi_c$:
\[
\dot V=e_{py}\dot e_{py}
=V_g e_{py}\sin(\chi-\chi_q)\cos\gamma
=-V_g e_{py}\sin\!\left(\chi_\infty\frac{2}{\pi}\tan^{-1}(k_{path}e_{py})\right)\cos\gamma<0.
\]
Hence $e_{py}\to 0$ asymptotically.

\subsection{Orbit path representation and kinematics}
Orbit:
\[
\mathcal{P}_{orbit}(c,\rho,\lambda)=\{r\in\R^3:\; r=c+\lambda\rho(\cos\varphi,\sin\varphi,0)^T,\;\varphi\in[0,2\pi)\}.
\]
Here:
\begin{itemize}
  \item $c$ = center, $\rho$ = radius,
  \item $\lambda=+1$ clockwise, $\lambda=-1$ counterclockwise,
  \item $d=\sqrt{(p_n-c_n)^2+(p_e-c_e)^2}$.
\end{itemize}
Polar kinematics:
\[
\dot d=V_g\cos(\chi-\varphi),\qquad
\dot\varphi=\frac{V_g}{d}\sin(\chi-\varphi).
\]
Desired tangent course:
\[
\chi_o=\varphi+\lambda\frac{\pi}{2}.
\]

\subsection{Orbit vector-field command and stability}
Course field:
\[
\chi_d(d-\rho,\lambda)=\chi_o+\lambda\tan^{-1}\!\left(k_{orbit}\frac{d-\rho}{\rho}\right),\quad k_{orbit}>0.
\]
Lyapunov candidate:
\[
V=\frac{1}{2}(d-\rho)^2.
\]
With $\chi=\chi_d$:
\[
\dot V=(d-\rho)V_g\cos(\chi_d-\varphi)
=-V_g(d-\rho)\sin\!\left(\tan^{-1}\!\left(k_{orbit}\frac{d-\rho}{\rho}\right)\right)<0,\quad d\neq \rho.
\]
Therefore $d\to \rho$ asymptotically.

\subsection{Practical angle wrapping issue}
To avoid command jumps near $\pm\pi$, choose integer wrapping index $m$ such that angle differences stay in $[-\pi,\pi]$.
For orbit phase:
\[
\varphi=\tan^{-1}\!\left(\frac{p_e-c_e}{p_n-c_n}\right)+2\pi m,\quad m\in\mathbb{Z}.
\]

\subsection{Classification of path-following methods}
\begin{itemize}
  \item Geometric methods: pure pursuit, LOS, NLGL.
  \item Dubins-curve-based methods.
  \item Vector-field-based methods.
\end{itemize}
Many methods rely on a \emph{virtual target point (VTP)} ahead on path.

\subsection{Carrot-chasing algorithm}
\textbf{Straight-line mode:}
\begin{itemize}
  \item project UAV position on line segment,
  \item place VTP at look-ahead distance $\delta$ along line,
  \item command heading toward VTP.
\end{itemize}
\textbf{Loiter mode:}
\begin{itemize}
  \item for desired circle of radius $r$, set VTP at angular offset $\lambda$ ahead on orbit,
  \item command heading toward VTP.
\end{itemize}

\subsection{Nonlinear Guidance Law (NLGL) / L1 Logic}
Look-ahead point at distance $L_1$, angle $\eta$ between velocity vector and line to look-ahead point:
\[
L_1=2R\sin\eta \;\Rightarrow\; R=\frac{L_1}{2\sin\eta}.
\]
Lateral acceleration command:
\[
a_{s,cmd}=\frac{V^2}{R}=\frac{2V^2}{L_1}\sin\eta.
\]

Small-angle perturbation about straight-line gives PD-like form:
\[
\eta\approx \frac{d-d^*_{ref.pt}}{L_1}+\frac{\dot d}{V},
\]
\[
a_{s,cmd}\approx \frac{2V}{L_1}\dot d+\frac{2V^2}{L_1^2}(d-d^*_{ref.pt}).
\]
Using $\ddot d=-a_{s,cmd}$:
\[
\ddot d+\frac{2V}{L_1}\dot d+\frac{2V^2}{L_1^2}d=\frac{2V^2}{L_1^2}d^*_{ref.pt}.
\]
Transfer function:
\[
\frac{d(s)}{d^*_{ref.pt}(s)}
=\frac{2(V/L_1)^2}{s^2+(2V/L_1)s+2(V/L_1)^2}.
\]
Hence
\[
\zeta=0.707,\qquad \omega_n=\frac{\sqrt{2}V}{L_1}.
\]

\subsection{Look-ahead distance selection insight}
\begin{itemize}
  \item Small $L_1$ gives high gain, faster response, potentially aggressive control.
  \item For sinusoidal path with wavelength $L_p$, practical recommendation: $L_1\lesssim L_p/4.4$ for good phase recovery near bandwidth.
\end{itemize}

\subsection{PD/PID vs NLGL in wind}
Observed in lecture comparisons:
\begin{itemize}
  \item fixed-gain PD/PID on cross-track error can show larger deviations in strong wind and varying ground speed;
  \item NLGL naturally uses ground speed in acceleration command, improving performance in upwind/downwind sectors.
\end{itemize}

\newpage
\section{Lecture 15: Circular NLGL Analysis, Path Manager, Fillets, Dubins Basics}

\subsection{Circular path with nonlinear guidance assumptions}
\begin{itemize}
  \item Inner-loop control is sufficiently fast.
  \item All angular quantities are signed consistently.
  \item Look-ahead distance $L_1$ is constant.
\end{itemize}
Key geometry:
\[
V_L=V\cos\eta,\qquad V_1=V\sin\eta,\qquad
V_T=\frac{V\cos\eta}{\cos\beta},\qquad V_2=V\cos\eta\tan\beta.
\]

\subsection{Dynamics of Eta and Beta}
\[
\dot\eta=-\frac{V}{L_1}\left(\sin\eta+\cos\eta\tan\beta\right),
\]
\[
\dot\beta=\frac{V}{L_1}\left(\sin\eta-\cos\eta\tan\beta\right)-\frac{V\cos\eta}{R\cos\beta}.
\]
Stationary point from $\dot\eta=\dot\beta=0$:
\[
\eta_0=-\beta_0,\qquad
\eta_0=\sin^{-1}\!\left(\frac{L_1}{2R}\right),\qquad
L_1\le 2R.
\]

\subsection{Linearization near circular equilibrium}
Perturbations:
\[
\eta=\eta_0+\Delta\eta,\qquad \beta=\beta_0+\Delta\beta.
\]
Linearized second-order behavior yields:
\[
\omega_n=\frac{\sqrt{2}V}{L_1},
\qquad
\zeta=\frac{1}{\sqrt{2}}\sqrt{1-\left(\frac{L_1}{2R}\right)^2}.
\]
Compared to straight-line case, $\omega_n$ stays same while $\zeta$ depends on $L_1/R$.

\subsection{Waypoint path manager motivation}
A waypoint list:
\[
\mathcal{W}=\{w_1,w_2,\dots,w_N\},\qquad w_i\in\R^3.
\]
Simple ``ball around waypoint'' switching can fail under wind and short segments.

\subsection{Half-plane switching}
Half-plane:
\[
\mathcal{H}(r,n)=\{p\in\R^3:(p-r)^T n\ge 0\}.
\]
Direction for segment $w_i\to w_{i+1}$:
\[
q_i=\frac{w_{i+1}-w_i}{\|w_{i+1}-w_i\|}.
\]
Switch at waypoint $w_i$ using:
\[
n_i=\frac{q_{i-1}+q_i}{\|q_{i-1}+q_i\|},
\]
and transition once UAV enters $\mathcal{H}(w_i,n_i)$.

\subsection{Fillet maneuver for smooth line transitions}
Turn angle:
\[
\varrho=\cos^{-1}(-q_{i-1}^Tq_i).
\]
For fillet radius $R$:
\[
c=w_i-\frac{R}{\sin(\varrho/2)}\frac{q_{i-1}-q_i}{\|q_{i-1}-q_i\|},
\]
\[
r_1=w_i-\frac{R}{\tan(\varrho/2)}q_{i-1},\qquad
r_2=w_i+\frac{R}{\tan(\varrho/2)}q_i.
\]
Procedure:
\begin{enumerate}
  \item follow line $w_{i-1}\to w_i$ to half-plane at $r_1$,
  \item follow circular arc (fillet) of radius $R$ about center $c$,
  \item switch to line $w_i\to w_{i+1}$ at half-plane through $r_2$.
\end{enumerate}

\subsection{Path-length correction with fillets}
Nominal polyline length:
\[
|W|=\sum_{i=2}^{N}\|w_i-w_{i-1}\|.
\]
Fillet-corrected length:
\[
|W|_F
=|W|+\sum_{i=2}^{N}\left[
R(\pi-\varrho_i)-\frac{2R}{\tan(\varrho_i/2)}
\right].
\]
This matters in cooperative timing and arrival synchronization.

\subsection{Dubins vehicle and path decomposition}
Bounded-curvature kinematics:
\[
\dot p_n=V\cos\vartheta,\qquad
\dot p_e=V\sin\vartheta,\qquad
\dot\vartheta=u,\quad u\in[-\bar u,\bar u].
\]
Minimum-turn radius:
\[
R=\frac{V}{\bar u}.
\]
Dubins path between $(p_s,\chi_s)$ and $(p_e,\chi_e)$ consists of two arcs and one straight segment (e.g., RSR, RSL, LSR, LSL).

\subsection{Circle centers and angle operator}
For pose $(p,\chi)$:
\[
c_r=p+R(\cos(\chi+\pi/2),\sin(\chi+\pi/2),0)^T,
\]
\[
c_l=p+R(\cos(\chi-\pi/2),\sin(\chi-\pi/2),0)^T.
\]
Use wrapped-angle operator:
\[
\langle\phi\rangle=\phi\;\text{mod}\;2\pi.
\]
CW and CCW arc angle measures are defined with this wrapping.

\subsection{Length formulas introduced in Lecture 15}
\textbf{RSR:}
\[
L_{RSR}=\|c_{rs}-c_{re}\|+d_1+d_2,
\]
with
\[
d_1=R\langle 2\pi+\langle\vartheta-\pi/2\rangle-\langle\chi_s-\pi/2\rangle\rangle,
\]
\[
d_2=R\langle 2\pi+\langle\chi_e-\pi/2\rangle-\langle\vartheta-\pi/2\rangle\rangle.
\]

\textbf{RSL:}
Let $l=\|c_{le}-c_{rs}\|$, and
\[
\vartheta_2=\vartheta-\left(\frac{\pi}{2}-\sin^{-1}\frac{2R}{l}\right).
\]
Then
\[
L_{RSL}=\sqrt{l^2-4R^2}+d_3+d_4,
\]
\[
d_3=R\langle 2\pi+\langle\vartheta_2\rangle-\langle\chi_s-\pi/2\rangle\rangle,
\]
\[
d_4=R\langle 2\pi+\langle\vartheta_2+\pi\rangle-\langle\chi_e+\pi/2\rangle\rangle.
\]

\newpage
\section{Lecture 16: Dubins Completion and UAV Path Planning Algorithms}

\subsection{Remaining Dubins cases}
\textbf{LSR:}
\[
L_{LSR}=\sqrt{l^2-4R^2}+d_5+d_6,
\]
\[
d_5=R\langle2\pi+\langle\chi_s+\pi/2\rangle-\langle\vartheta+\vartheta_2\rangle\rangle,
\]
\[
d_6=R\langle2\pi+\langle\chi_e-\pi/2\rangle-\langle\vartheta+\vartheta_2-\pi\rangle\rangle.
\]

\textbf{LSL:}
\[
L_{LSL}=\|c_{ls}-c_{le}\|+d_7+d_8,
\]
\[
d_7=R\langle2\pi+\langle\chi_s+\pi/2\rangle-\langle\vartheta+\pi/2\rangle\rangle,
\]
\[
d_8=R\langle2\pi+\langle\vartheta+\pi/2\rangle-\langle\chi_e+\pi/2\rangle\rangle.
\]
Choose minimum among feasible RSR/RSL/LSR/LSL.

\subsection{Path planning classes}
\begin{itemize}
  \item \textbf{Deliberative planning}: uses global map/model; explicit path computation; good for structured known environments.
  \item \textbf{Reactive planning}: uses local sensing and behavior rules; better for uncertain dynamic environments.
\end{itemize}
For many fixed-wing UAV missions, efficient deliberative planning with simplified motion/wind models is preferred.

\subsection{Point-to-point planning with Voronoi graph}
Given obstacle point set $Q$, Voronoi partition creates cells and edges equidistant from obstacles, producing safer corridors.
Graph:
\[
G=(V,E),\quad V^+=V\cup\{p_s,p_e\}.
\]
Start/end nodes are connected to nearest graph nodes.

\subsection{Edge cost with safety-distance penalty}
For edge $(v_1,v_2)$ and obstacle point $p$, parameterize segment:
\[
w(\sigma)=(1-\sigma)v_1+\sigma v_2,\quad \sigma\in[0,1].
\]
Distance minimization gives
\[
D=\min_{\sigma\in[0,1]}\|p-w(\sigma)\|.
\]
Unconstrained optimizer:
\[
\sigma^*=-\frac{(p-v_1)^T(v_1-v_2)}{\|v_1-v_2\|^2}.
\]
Use clipped distance $D'(v_1,v_2,p)$:
\begin{itemize}
  \item $D(\sigma^*)$ if $\sigma^*\in[0,1]$,
  \item $\|p-v_1\|$ if $\sigma^*<0$,
  \item $\|p-v_2\|$ if $\sigma^*>1$.
\end{itemize}
Distance from edge to obstacle set:
\[
D(v_1,v_2,Q)=\min_{p\in Q}D'(v_1,v_2,p).
\]
Edge cost:
\[
J(v_1,v_2)=k_1\|v_1-v_2\|+\frac{k_2}{D(v_1,v_2,Q)},\quad k_1,k_2>0.
\]

\subsection{Dijkstra search procedure}
\begin{enumerate}
  \item Initialize tentative distance 0 at start node, $\infty$ for others.
  \item Repeatedly select unvisited node with minimum tentative distance.
  \item Relax all unvisited neighbors.
  \item Mark node visited.
  \item Stop when destination visited or remaining tentative distances are $\infty$.
\end{enumerate}

\subsection{RRT for point-to-point planning}
Rapidly Exploring Random Tree:
\begin{itemize}
  \item randomly sample a configuration in workspace,
  \item find nearest node in existing tree,
  \item extend by fixed step $D$ toward sample,
  \item add segment only if collision-free,
  \item continue until close enough to goal.
\end{itemize}
RRT extends naturally to nonlinear and higher-dimensional models, including 3D with climb/descent constraints.

\subsection{Coverage path planning concept}
Goal is area coverage (search), not just start-to-goal transit.
\begin{itemize}
  \item Maintain terrain map (collision check) and return map (search benefit).
  \item Update return value after visit:
  \[
  \Upsilon_i[k]=\Upsilon_i[k-1]-c,\quad c>0.
  \]
  \item Build finite-depth look-ahead trees using motion primitives:
  step length $L$ and heading changes $\{0,\pm\nu\}$ (example: $\nu=\pi/6$).
\end{itemize}

\newpage
\section{Lecture 17: Vision-Guided Navigation, Geolocation, and Time-to-Collision}

\subsection{Role of onboard vision}
EO/IR cameras support:
\begin{itemize}
  \item surveillance and situational awareness,
  \item guidance/control by using image features as measurements.
\end{itemize}

\subsection{Frames and transformation chain}
\begin{itemize}
  \item Body frame: $F_b$
  \item Gimbal frame: $F_g$
  \item Camera frame: $F_c$
\end{itemize}
Body-to-gimbal-1 (azimuth rotation $\alpha_{az}$):
\[
R_b^{g1}(\alpha_{az})=
\begin{bmatrix}
\cos\alpha_{az} & \sin\alpha_{az} & 0\\
-\sin\alpha_{az} & \cos\alpha_{az} & 0\\
0 & 0 & 1
\end{bmatrix}.
\]
Gimbal-1 to gimbal (elevation $\alpha_{el}$):
\[
R_{g1}^{g}(\alpha_{el})=
\begin{bmatrix}
\cos\alpha_{el} & 0 & -\sin\alpha_{el}\\
0 & 1 & 0\\
\sin\alpha_{el} & 0 & \cos\alpha_{el}
\end{bmatrix}.
\]
Hence
\[
R_b^g=R_{g1}^g(\alpha_{el})R_b^{g1}(\alpha_{az}).
\]
Camera alignment from gimbal basis:
\[
i_c=j_g,\quad j_c=k_g,\quad k_c=i_g
\Rightarrow
R_g^c=
\begin{bmatrix}
0&1&0\\
0&0&1\\
1&0&0
\end{bmatrix}.
\]

\subsection{Camera projection model}
For image width $M$ pixels and field of view $\nu$:
\[
f=\frac{M}{2\tan(\nu/2)}.
\]
For pixel location $(\epsilon_x,\epsilon_y)$:
\[
F=\sqrt{f^2+\epsilon_x^2+\epsilon_y^2},
\qquad
l_c=\frac{L}{F}\begin{bmatrix}\epsilon_x\\ \epsilon_y\\ f\end{bmatrix},
\qquad
\check l_c=\frac{1}{F}\begin{bmatrix}\epsilon_x\\ \epsilon_y\\ f\end{bmatrix}.
\]
Monocular camera gives direction $\check l_c$, not absolute distance $L$.

\subsection{Gimbal pointing control}
Pan-tilt kinematics:
\[
\dot\alpha_{az}=u_{az},\qquad \dot\alpha_{el}=u_{el}.
\]
From desired LOS direction in body frame $\check l_d^b=[\check l_{xd}^b,\check l_{yd}^b,\check l_{zd}^b]^T$:
\[
\alpha_{az}^c=\tan^{-1}\!\left(\frac{\check l_{yd}^b}{\check l_{xd}^b}\right),\qquad
\alpha_{el}^c=-\sin^{-1}(\check l_{zd}^b).
\]
Simple proportional servo commands:
\[
u_{az}=k_{az}(\alpha_{az}^c-\alpha_{az}),\qquad
u_{el}=k_{el}(\alpha_{el}^c-\alpha_{el}).
\]

\subsection{Geolocation setup}
Object position:
\[
p_{obj}^i=p_{MAV}^i+R_b^iR_g^bR_c^g\,l_c
=p_{MAV}^i+L\,R_b^iR_g^bR_c^g\,\check l_c.
\]
Unknown $L$ causes scale ambiguity.

\subsection{Flat-earth geolocation estimate}
If camera-ground intersection uses known vehicle height $h$ and flat-earth assumption:
\[
L=\frac{h}{\cos\phi},
\qquad
\cos\phi=k_i^T\left(R_b^iR_g^bR_c^g\check l_c\right).
\]
Then:
\[
p_{obj}^i
=
\begin{bmatrix}p_n\\p_e\\p_d\end{bmatrix}
+
h\,
\frac{R_b^iR_g^bR_c^g\check l_c}{k_i^T\left(R_b^iR_g^bR_c^g\check l_c\right)}.
\]

\subsection{EKF formulation for geolocation}
State:
\[
\hat x=\begin{bmatrix}\hat p_{obj}^i\\ \hat L\end{bmatrix}.
\]
Prediction equations (stationary object assumption):
\[
\dot{\hat p}_{obj}^i=0,\qquad
\dot{\hat L}
=-\frac{(\hat p_{obj}^i-\hat p_{MAV}^i)^T\dot{\hat p}_{MAV}^i}{\hat L}.
\]
Measurement relation:
\[
p_{MAV}^i=p_{obj}^i-LR_b^iR_g^bR_c^g\check l_c.
\]
Jacobian terms follow from these dynamics and measurement model.

\subsection{Target pixel filtering and velocity estimation}
Raw pixel vector $\bar\epsilon=[\bar\epsilon_x,\bar\epsilon_y]^T$ is noisy.
Use low-pass and filtered derivative:
\[
\epsilon(s)=\frac{1}{\tau s+1}\bar\epsilon,\qquad
\dot\epsilon(s)=\frac{s}{\tau s+1}\bar\epsilon.
\]
Using Tustin discretization:
\[
\epsilon[n]=\frac{2\tau-T_s}{2\tau+T_s}\epsilon[n-1]
+\frac{T_s}{2\tau+T_s}\left(\bar\epsilon[n]+\bar\epsilon[n-1]\right),
\]
\[
\dot\epsilon[n]=\frac{2\tau-T_s}{2\tau+T_s}\dot\epsilon[n-1]
+\frac{2}{2\tau+T_s}\left(\bar\epsilon[n]-\bar\epsilon[n-1]\right).
\]

\subsection{Removing apparent motion due to platform/gimbal rotation}
Rate transport:
\[
\frac{d\check l_c}{dt_i}
=\frac{d\check l_c}{dt_c}+\omega_{c/i}^c\times \check l_c.
\]
Normalized LOS derivative from pixel rates:
\[
\frac{d\check l_c}{dt_c}=Z(\epsilon)\dot\epsilon,
\]
where
\[
Z(\epsilon)=\frac{1}{F^3}
\begin{bmatrix}
\epsilon_y^2+f^2 & -\epsilon_x\epsilon_y\\
-\epsilon_x\epsilon_y & \epsilon_x^2+f^2\\
-\epsilon_x f & -\epsilon_y f
\end{bmatrix}.
\]
Subtracting rotation-induced term isolates translationally useful motion.

\subsection{Time-to-collision (TTC)}
Definition:
\[
t_c=\frac{L}{\dot L}.
\]
With monocular camera only, TTC is ambiguous unless additional information is used.

If object physical size $S_{obj}$ and image size $\epsilon_s$ are available:
\[
\frac{S_{obj}}{L}=\frac{\epsilon_s}{F}.
\]
Differentiation gives:
\[
\frac{\dot L}{L}
=\frac{\epsilon_x\dot\epsilon_x+\epsilon_y\dot\epsilon_y}{F}
-\frac{\dot\epsilon_s}{\epsilon_s},
\]
hence
\[
t_c
=\frac{1}{\dfrac{\epsilon_x\dot\epsilon_x+\epsilon_y\dot\epsilon_y}{F}
-\dfrac{\dot\epsilon_s}{\epsilon_s}}.
\]

\newpage
\section{Consolidated Revision Checklist (Lectures 13--17)}
\begin{enumerate}
  \item Derive and interpret coordinated-turn and accelerated-climb equations.
  \item Explain why path following is preferred over trajectory tracking for small UAVs in wind.
  \item Reproduce straight-line and orbit vector-field commands and Lyapunov proofs.
  \item Derive NLGL linearized transfer function and identify $\zeta,\omega_n$.
  \item Compute half-plane switches and fillet geometry from waypoint triplets.
  \item Evaluate Dubins path length for all four cases and pick minimum feasible route.
  \item Build Voronoi graph edge costs and describe Dijkstra steps.
  \item Explain RRT growth, collision checking, and stopping criteria.
  \item Write complete frame chain for vision geolocation and identify source of scale ambiguity.
  \item Derive TTC from pixel-size change and discuss assumptions.
\end{enumerate}

\section*{Primary Reference}
Randal Beard and Timothy W. McLain, \emph{Small Unmanned Aircraft: Theory and Practice}, Princeton University Press, 2012.

\end{document}
